\documentclass[a4paper,12pt]{article} 

%%%%%%%%%%%%%%%%%%%%%%%%%%%%%%%% CONSTANTES %%%%%%%%%%%%%%%%%%%%%%%%%%%%%%%%%%%
\newcommand{\numero}{1}                                    %Numéro de la série -1

\usepackage[french]{babel}
\usepackage[utf8]{inputenc}
\usepackage{answers}

\usepackage{hyperref}
\usepackage{multicol}

\usepackage[table,xcdraw]{xcolor}
\usepackage{listings}
\definecolor{ForestGreen}{RGB}{34,139,34}


\usepackage{enumitem}

\AtBeginDocument{\def\labelitemi{$\bullet$}}


\newcommand{\py}{\lstinline{Python} }


\definecolor{backcolour}{rgb}{0.95,0.95,0.92}

\lstset{%
	language         = Python,
	backgroundcolor  = \color{backcolour},
	basicstyle       = \ttfamily, % \upshape\ttfamily,
	keywordstyle     = \bfseries\color{blue}, %\bfseries,
	stringstyle      = \color{magenta},
	commentstyle     = \color{ForestGreen},
	alsoletter = > ,
	morekeywords = {>>>,as,assert,False,None, nonlocal,True, with,yield , <<, >>, :},
	showstringspaces = false,
	numbers=left,
	stepnumber=1,
	literate={à}{{\`{a}}}1 {é}{{\'e}}1 {è}{{\`{e}}}1 {ê}{{\^{e}}}1 {Ê}{{\^{E}}}1 {î}{{\^i}}1 {ô}{{\^{o}}}1 {ç}{{\c{c}}}1 {Ç}{{\c{C}}}1
}

\newcommand{\itemb}[1]{\item \textbf{#1}}

\usepackage{fancyhdr}  %package pour en-tetes et pied de pages
\usepackage{sectsty} % Permet de faire des modifications de police dans diverses sections des "headings" (cf. modif presentation de la page)
\pagestyle{fancy}       %Style pour en-tetes et pieds de pages
\fancyhead[CO,CE]{\sc Série 1\hspace{0.5mm}}
\fancyhead[RO,LE]{Collège Sismondi}  % LaTeX/TEX define \strut to be an invisible box of width zero that extends just enough above and below the baseline. Cela permet d'augementer légèrement la taille en bas de la box de manière à ce qu'elle soit collée à la ligne.
\fancyhead[LO,RE]{\small\ \textsl{1\textsuperscript{ère} année - DO Informatique}}
\fancyfoot[RO,LE]{2021 - 2022}
\fancyfoot[LO,RE]{\small }
\fancyfoot[CO,CE]{\thepage}

\fancyhfoffset[l]{1.2cm} % le "l" en paramètre permet d'indiquer qu'on ne veut modifier que la marge à gauche.
\renewcommand{\headrule}{{%
		\hrule \headwidth \headrulewidth \vskip-\headrulewidth}}
\renewcommand\footrulewidth{\headrulewidth}
\renewcommand{\footrule}{{%
		\vskip-\footruleskip\vskip-\footrulewidth
		\hrule \headwidth \footrulewidth\vskip\footruleskip}}

\usepackage{tikz}
%-------------------------------------------------------------------------------
%---- Eclairage : en encadré sur fond jaune avec symbôle "ampoule" à gauche ----
%-------------------------------------------------------------------------------
\definecolor{coleclairage}{RGB}{255 , 221 , 156}
\definecolor{contoureclairage}{RGB}{255 , 192 , 0}
\newenvironment{eclairage}
{
	\begin{center}%
		\begin{tikzpicture}%
			\node[rectangle, draw=contoureclairage, top color=coleclairage!50, bottom color=coleclairage!140, rounded corners=5pt, inner xsep=5pt, inner ysep=6pt, outer ysep=10pt]\bgroup                     
			\begin{minipage}{0.98\linewidth}
				\begin{minipage}{0.08\linewidth}\centerline{\includegraphics[scale=1]{Symbole_eclairage.png}}\end{minipage}
				\begin{minipage}{0.89\linewidth}\itshape\footnotesize
				}
				{                		
				\end{minipage}
			\end{minipage}\egroup;%
		\end{tikzpicture}%
	\end{center}%
}

%-------------------------------------------------------------------------------
%---- apprendre : en encadré sur fond jaune avec symbôle "ampoule" à gauche ----
%-------------------------------------------------------------------------------
\definecolor{colapprendre}{RGB}{50,205,50}
\definecolor{contourapprendre}{RGB}{34,139,34}
\newenvironment{apprendre}
{
	\begin{center}%
		\begin{tikzpicture}%
			\node[rectangle, draw=contourapprendre, top color=colapprendre!10, bottom color=colapprendre!50, rounded corners=5pt, inner xsep=5pt, inner ysep=6pt, outer ysep=10pt]\bgroup                     
			\begin{minipage}{0.98\linewidth}
				\begin{minipage}{0.08\linewidth}\centerline{\includegraphics[width=30px]{Symbole_learn.png}}\end{minipage}
				\begin{minipage}{0.89\linewidth}\itshape\footnotesize
				}
				{                		
				\end{minipage}
			\end{minipage}\egroup;%
		\end{tikzpicture}%
	\end{center}%
}

\definecolor{colimportant}{RGB}{247 , 189 , 164}
\definecolor{contourimportant}{RGB}{237 , 125 , 49}
\newenvironment{important}
{
	\begin{center}%
		\begin{tikzpicture}%
			\node[rectangle, draw=contourimportant, top color=colimportant!50, bottom color=colimportant!140, rounded corners=5pt, inner xsep=5pt, inner ysep=6pt, outer ysep=10pt]\bgroup                     
			\begin{minipage}{0.08\linewidth}\centerline{\includegraphics[scale=0.8]{Symbole_attention.png}}\end{minipage}
			\begin{minipage}{0.89\linewidth}
			}
			{                		
			\end{minipage}\egroup;
		\end{tikzpicture}%
	\end{center}%
}

%-----------------------------------------------------------------
%---- Modification présentation de la page: marges de la page ----
%-----------------------------------------------------------------
%\addtolength{\hoffset}{-1in}              % 1
%\addtolength{\voffset}{-1in}              % 2
\addtolength{\oddsidemargin}{-0.1 in} % 3
\addtolength{\evensidemargin}{-1in} % 3
\addtolength{\topmargin}{-1in}       % 4
\addtolength{\headheight}{6pt}       % 5
%\addtolength{\headsep}{-0.2cm}           % 6
\setlength{\textheight}{26cm}    % 7
\setlength{\textwidth}{16.5cm}      % 8
\addtolength{\marginparsep}{0pt}      % 9
\setlength{\marginparwidth}{0pt}   % 10
\addtolength{\footskip}{-1mm}           %11

\setlength{\parindent}{0em}% pas d'indentation


% Customiser le nom des sections
\usepackage{titlesec}
\titleformat{\section}[hang]{\Large \bfseries}{Série \thesection:\ }{0pt}{}

\renewcommand{\familydefault}{\sfdefault} % pour avoir des polices san serif

\newtheorem{Exc}{Exercice}
\Newassociation{correction}{Soln}{mycor}
\renewcommand{\Solnlabel}[1]{\bfseries Ex #1 }
\def\exo#1{%
	\futurelet\testchar\MaybeOptArgmyexoo}
\def\MaybeOptArgmyexoo{
	\ifx[\testchar \let\next\OptArgmyexoo
	\else \let\next\NoOptArgmyexoo \fi \next}
\def\OptArgmyexoo[#1]{%
	\begin{Exc}[#1]\normalfont}
	\def\NoOptArgmyexoo{%
		\begin{Exc}\normalfont}
		\newcommand{\finexo}{\end{Exc} \vspace{3mm}}
	\newcommand{\flag}[1]{}
	\newcommand{\entete}[1]

\newcommand{\getexocompteur}{{\the\numexpr \arabic{Exc}  \relax}}	
	
\newcommand{\eexo}{\vspace{5mm}} % espace pour séparer les exercices


\begin{document}
%		\title{\vspace{-3cm}Série 1}
%		\date{\vspace{-2cm}}
%		\maketitle

\fancyhead[CO,CE]{\sc Série \arabic{section} \hspace{0.5mm}}

\setcounter{section}{\numero}
\section{Types et opérations}				
\Opensolutionfile{mycor}[cor01]

%%%%%%%%%%%%%%%%%%%%%%%%%%%%%%%% E X E R C I C E  %%%%%%%%%%%%%%%%%%%%%%%%%%%%%%%%%%%
\exo{}  ~\\ 
Répondre aux questions suivantes et ensuite les vérifier dans la console Python.

\begin{enumerate}
	\item De quel type est le résultat du calcul 72/3 en Python ?
	\item Quelle est la valeur des expressions suivantes:  
	   \begin{lstlisting}[numbers=none]
>>> (2 or 0) and (0.1 == 1/10)
>>> (2 or 0) and (0.1 == (0.4-0.3))
>>> False or not (not True)
       \end{lstlisting}
       Pouvez vous expliquer la différence entre la première et la deuxième expression ?
	\end{enumerate}

	\begin{correction}
		~\\ 
		\begin{enumerate}
	\item \lstinline{float} même si 72 et 3 sont du type \lstinline{int}. Pour vérifier: \lstinline{>>> type(72/3)}
	\item \lstinline{>>> True}\newline
	      \lstinline{>>> False}\newline
	      \lstinline{>>> True}\newline
	      Le problème est que le nombre \lstinline{0,4} est une approximation et sa soustraction avec \lstinline{0,3} donne une valeur \lstinline{> 0,1}.
 		\end{enumerate}
	\end{correction}
\finexo
%%%%%%%%%%%%%%%%%%%%%%%%%%%%%%%% E X E R C I C E %%%%%%%%%%%%%%%%%%%%%%%%%%%%%%%%%%%
\exo{}  ~\\ 
Tester dans la console \py les commandes suivantes:
\begin{lstlisting}[numbers=none]
>>> type(123)
>>> type(123.)
>>> type("blabla")
>>> type(True)
\end{lstlisting}

Indiquer lesquels des objets suivants sont valides en Python et le cas échéant de quel type d'objet il s'agit :
\begin{multicols}{2}
	\begin{enumerate}[label=\alph*)]
		\item "rewr"
		\item gdru
		\item 5
		\item 5.0
		\item 5,0
		\item 'sism2021'
		\item 18gd
		\item "65.5"
		\item True
		\item "False"
	\end{enumerate}
\end{multicols}
	\begin{correction}
	~\\ 
	Utiliser la fonction \lstinline{type()} dans l'interpréteur \py pour vérifier vos réponses.
	\begin{multicols}{2}
	\begin{enumerate}[label=\alph*)]
		\item "rewr": str
		\item Erreur grdu n'est pas défini.
		\item 5: int
		\item 5.0: float
		\item 5,0: n'est pas un float mais un tuple (sera étudié plus tard)
		\item 'sism2021': est un str
		\item 18gd: invalide syntaxe
		\item "65.5": str
		\item True: : bool (booléen)
		\item "False": str
	\end{enumerate}
\end{multicols}
\end{correction}
\finexo

%%%%%%%%%%%%%%%%%%%%%%%%%%%%%%%% E X E R C I C E %%%%%%%%%%%%%%%%%%%%%%%%%%%%%%%%%%%
\exo{}[Convertir dans un type]  ~\\ 
 En s'aidant de la console \py, essayez de déterminer le rôle des fonctions \lstinline{int()},
\lstinline{float()} et \lstinline{str()} en écrivant par exemple les lignes suivantes :
\begin{lstlisting}[numbers=none]
>>> int(3.6)
>>> float(3)
>>> str(3)
>>> 2**1000
>>> float(2**1000)
>>> int('3')
\end{lstlisting}
\vspace{0.3cm}
\textit{Indication : Vous pouvez vous aider de la commande \lstinline{type()} pour analyser le résultat des fonctions \lstinline{int()}, \lstinline{float()} et \lstinline{str()}.}
	\begin{correction}
	~\\ 
	Les fonctions \lstinline{int()}, \lstinline{float()} et \lstinline{str()} permettent de convertir les données dans le type indiqué par la fonction.
\end{correction}
\finexo
\newpage
%%%%%%%%%%%%%%%%%%%%%%%%%%%%%%%% E X E R C I C E %%%%%%%%%%%%%%%%%%%%%%%%%%%%%%%%%%%
\exo{}[Expression et priorité des opérations]  ~\\ 
Prédire le résultat des instructions suivantes en \py:
\begin{multicols}{2}
	\begin{enumerate}[label=\alph*)]
		\item \lstinline{3+2*5}
		\item \lstinline{2*5**2}
		\item \lstinline{(2*5)**2}
		\item \lstinline{6/4*2}
		\item \lstinline{6.0/4*2}
		\item \lstinline{7 + 3 * 6 / 2 - 1}
		\item \lstinline{2 % 2 + 2 * 2 - 2 / 2}
		\item \lstinline{3 * 9 * (3 + (9 * 3 / (3)))}
	\end{enumerate}
\end{multicols}
	Entrer ces expressions dans la console \py pour vérifier les réponses.
	
    \begin{correction}
	~\\ 
\begin{multicols}{2}
	\begin{enumerate}[label=\alph*)]
		\item \lstinline{13}
		\item \lstinline{50}
		\item \lstinline{100}
		\item \lstinline{3.0}
		\item \lstinline{3.0}
		\item \lstinline{15.0}
		\item \lstinline{3.0}
		\item \lstinline{324.0}
		\end{enumerate}
	\end{multicols}
\end{correction}
\finexo

%%%%%%%%%%%%%%%%%%%%%%%%%%%%%%%% E X E R C I C E %%%%%%%%%%%%%%%%%%%%%%%%%%%%%%%%%%%
\exo{}  ~\\ 
Parmi les propositions, laquelle n'est pas une expression ?
\begin{multicols}{2}
	\begin{enumerate}[label=\alph*)]
		\item \lstinline{a < b}
		\item \lstinline{a != b}
		\item \lstinline{a = b}
		\item \lstinline{a >= b}
	\end{enumerate}
\end{multicols}
\begin{correction}
	~\\ 
	c) a = b n'est pas une expression, il s'agit d'une affectation. La valeur de b est affectée à la variable a.
\end{correction}
\finexo
%%%%%%%%%%%%%%%%%%%%%%%%%%%%%%%% E X E R C I C E %%%%%%%%%%%%%%%%%%%%%%%%%%%%%%%%%%%

\exo{}  ~\\ 
On considère les instructions suivantes exécutées dans l'ordre.
Quel est le résultat affiché de la commande: \lstinline{print(a==b+1)}?
\begin{multicols}{2}
	\begin{enumerate}[label=\alph*)]
		\item \lstinline{8}
		\item \lstinline{Aucun, une erreur est signalée.}
		\item \lstinline{False}
		\item \lstinline{True}
	\end{enumerate}
\end{multicols}
\begin{lstlisting}[numbers=none]
>>> a=8
>>> b=5
>>> a==b+1
>>> b=b+1
>>> a==b+1
>>> b=b+1
>>> print(a==b+1)
\end{lstlisting}
\begin{correction}
	~\\ 
	d) \lstinline{True}
\end{correction}
\finexo
%%%%%%%%%%%%%%%%%%%%%%%%%%%%%%%% E X E R C I C E %%%%%%%%%%%%%%%%%%%%%%%%%%%%%%%%%%%
\exo{}  ~\\ 
On considère les deux instructions a=a+b et b=a-b exécutées dans l'ordre. Quelle affirmation est exacte ?
	\begin{enumerate}[label=\alph*)]
		\item Si les valeurs initiales de a et b son respectivement "bon" et "jour", alors le programme est interrompu par une erreur.
		\item Si les valeurs initiales de a et b son respectivement "bon" et "jour", alors les valeurs finales (des variables a et b) sont respectivement "bonjour" et "bon".
		\item Si les valeurs initiales de a et b son respectivement 5 et 2, alors les valeurs finales (des variables a et b) sont respectivement 7 et 3.
		\item Si les valeurs initiales de a et b son respectivement 5 et 2, alors les valeurs finales (des variables a et b) sont respectivement 7 et 2.
	\end{enumerate}

\begin{correction}
	~\\ 
	L'affirmation a) est correct car la soustraction (-) ne fait pas partie des opérations permises avec des chaînes de caractères.
\end{correction}
\finexo
%%%%%%%%%%%%%%%%%%%%%%%%%%%%%%%% E X E R C I C E %%%%%%%%%%%%%%%%%%%%%%%%%%%%%%%%%%%
\exo{}  ~\\ 
Manipulations fondamentales sur les chaînes. On définit les deux suivantes:
\begin{lstlisting}[numbers=none]
>>> ch1 = "Ma première chaîne de caractères"
>>> ch2 = "en Python"
\end{lstlisting}
Répondez, uniquement à l'aide de \lstinline{ch1} et \lstinline{ch2}, aux questions suivantes. Donnez le code \py correspondant:
	\begin{enumerate}[label*=\arabic*] 
		\item Quel est le nombre de caractères de \lstinline{ch1} ?
		\item Créer la chaîne \lstinline{ch3} égale à "première".
		\item Créer la chaîne \lstinline{ch4} égale à "Ma première chaîne de caractères en Python".
		\item Créer la chaîne \lstinline{ch5} égale à "Mon premier".
	\end{enumerate}

\begin{correction}
	~\\ 
	\begin{enumerate}[label*=\arabic*] 
		\item \lstinline{>>> len(ch1)         # renvoie 32}
		\item \lstinline{>>> ch3 = ch1[3:11]}
		\item \lstinline{>>> ch4 = ch1 + ch1[2]+ ch2    # Ne pas oublier l'espace}
		\item 
		     \lstinline{>>> l2 = len(ch2)   # pour éviter de compter les indices dans ch2}\newline
		     \lstinline{>>> ch5 = ch1[0]+ch2[l2-2:]+ch1[2:8]+ch1[5]+ch1[4]}
	\end{enumerate}
\end{correction}
\finexo

% Solution		
		\newpage
		\setcounter{page}{1}
		\setcounter{section}{\numero}
		\Closesolutionfile{mycor}
		\titleformat{\section}[hang]{\Large \bfseries}{Corrigé Série \thesection:\ }{0pt}{}
		

		\fancyhead[CO,CE]{\sc Corrigé Série \arabic{section} \hspace{0.5mm}}
		\section{}
		\Readsolutionfile{mycor}
	\end{document}